% Volume VI: Adversaries, Platforms, and Automated Judgment
% Part XII. Platforms and Automated Publicity
% Chapter 65: Platforms as Evidentiary Institutions
% Spine: proof-spines/ch065-spine.md

\chapter{Platforms as Evidentiary Institutions}
\label{ch:ch065}


Platforms already function as \textbf{evidentiary institutions}. They collect records, rank claims, determine visibility, enforce categories, and preserve or erase histories long before any court, newsroom, or public inquiry confronts the material. What later counts as available evidence is therefore partly a consequence of platform architecture.

This role is quasi-legal even where it is not legally acknowledged. Decisions about retention, deletion, ranking, metadata, and categorical labeling shape chain of custody, salience, and retrievability, but the corresponding procedural safeguards are often thin or absent. The preface claim is accordingly that platforms do not merely host speech: they organize testimony, memory, and disappearance at scale, and so exercise evidentiary power that warrants institutional scrutiny.
